\subsection{APPROXIMATION OF CONTINUOUS FUNCTIONS BY POLYNOMIALS}

Given a set H of real-valued functions defined on a set X, we say that a real-valued function defined on X is a polynomial (resp. a polynomial with no constant term) with real coefficients, in the functions of H, if it is of the form \(x \to g(f_{1}(x), f_{2}(x), \ldots, f_{n}(x))\) where g is a polynomial (resp. a polynomial with no constant term) in n indeterminates (n arbitrary) with real coefficients, and the \(f_{i}\) ( \(1 \leqslant i \leqslant n\) ) belong to H.

THEOREM 3 (Weierstrass-Stone). Let X be a compact space and let H be a set of continuous real-valued functions on X which separates the points of X. Then every continuous real-valued function on X can be uniformly approximated by polynomials (with real coefficients) in the functions of H.

An equivalent statement of the theorem is that any subalgebra of \(\mathcal{C}(\mathbf{X};\mathbf{R})\) which contains the constant functions and separates the points of \(\mathbf{X}\) is dense in \(\mathcal{C}(\mathbf{X};\mathbf{R})\) .

Let \(H_{0}\) be the set of all polynomials in the functions of H, and let \(\overline{H}_{0}\) be the closure of \(H_{0}\) in C. If g is any polynomial in n variables with real coefficients, then \((u_{1}, u_{2}, \ldots, u_{n}) \to g(u_{1}, u_{2}, \ldots, u_{n})\) is a continuous mapping of \(C^{n}\) into C, which maps \(H_{0}^{n}\) into \(H_{0}\) , and therefore maps \(\overline{H}_{0}^{n}\) into \(\overline{H}_{0}\) (Chapter I, § 2, no. 1, Theorem 1). In particular, \(\overline{H}_{0}\) is a vector subspace of C and evidently satisfies the first and third conditions of Theorem 2; we shall show that it also satisfies the second condition, and this will prove that \(\overline{H}_{0} = C\) .

Since every function \(u \in \overline{\mathbf{H}}_0\) is bounded in \(\mathbf{X}\) , it is enough to prove the following lemma:

LEMMA 1. For each real number \(\varepsilon > 0\) and each compact interval \(\mathbf{I} \subset \mathbf{R}\) there exists a polynomial \(p(t)\) with no constant term such that \(|p(t) - |t|| \leqslant \varepsilon\) for all \(t \in \mathbf{I}\) .

It is enough to prove the lemma for an interval of the form \(\mathbf{I} = [-a, +a]\) and hence, replacing \(t\) by \(at\) , for the interval \(\mathbf{I} = [-1, +1]\) . Since \(|t| = \sqrt{t^2}\) , Lemma 1 is a consequence of the following result:

LEMMA 2. Let \((P_n)\) be the sequence of polynomials without constant terms defined by

\[
p _ {0} (t) = 0, \quad p _ {n + 1} (t) = p _ {n} (t) + \frac {1}{2} (t - (p _ {n} (t)) ^ {2}), \quad n \geqslant 0.\tag{1}
\]

In the interval {[}0, 1{]} the sequence \((p_n)\) is increasing and converges uniformly to \(\sqrt{t}\) .

To prove Lemma 2 it is enough to show that, for all \(t \in [0, 1]\) , we have

\[
0 \leqslant \sqrt {t} - p _ {n} (t) \leqslant \frac {2 \sqrt {t}}{2 + n \sqrt {t}},\tag{2}
\]

for (2) implies that \(0 \leqslant \sqrt{t} - p_n(t) \leqslant 2 / n\) .

We prove (2) by induction on \(n\) . It is true for \(n = 0\) . If \(n \geqslant 0\) it follows from the inductive hypothesis (2) that \(0 \leqslant \sqrt{t} - p_n(t) \leqslant \sqrt{t}\) , hence \(0 \leqslant p_n(t) \leqslant \sqrt{t}\) , and therefore from (1) we have

\[
\sqrt {t} - p _ {n + 1} (t) = (\sqrt {t} - p _ {n} (t)) \left(1 - \frac {1}{2} (\sqrt {t} + p _ {n} (t))\right),
\]

so that \(\sqrt{t} - p_{n+1}(t) \geqslant 0\) , and from (2)

\[
\begin{array}{r l} \sqrt {t} - p _ {n + 1} (t) & \leqslant \frac {2 \sqrt {t}}{2 + n \sqrt {t}} \left(1 - \frac {\sqrt {t}}{2}\right) \leqslant \frac {2 \sqrt {t}}{2 + n \sqrt {t}} \left(1 - \frac {\sqrt {t}}{2 + (n + 1) \sqrt {t}}\right) \\ & = \frac {2 \sqrt {t}}{2 + (n + 1) \sqrt {t}}. \end{array} \tag {Q.E.D.}
\]

If X is not compact, the conclusion of Theorem 3 is not necessarily valid. For example, a continuous real-valued function on R which is bounded and not constant cannot be uniformly approximated in R by polynomials (cf.~Exercise 6).

PROPOSITION 3. Let \((\mathbf{K}_t)_{t \in \mathbf{I}}\) be a family of compact intervals of \(\mathbf{R}\) , \(\mathbf{K} = \prod_{t \in \mathbf{I}} \mathbf{K}_t\) their product, and let \(\mathbf{X}\) be a compact subspace of \(\mathbf{K}\) . Then every continuous real-valued function on \(\mathbf{X}\) can be uniformly approximated by polynomials in the coordinates \(x_t = \mathrm{pr}_t x\) .

For if \(x = (x_{t})\) and \(y = (y_{t})\) are two distinct points of X, there is at least one index t such that \(x_{t} \neq y_{t}\) ; hence the family of continuous functions \(pr_{t}\) satisfies the conditions of Theorem 3.

PROPOSITION 4. Let X be a compact space, let A be a closed subspace of X, and let H be a set of continuous real-valued functions on X which separates the points of \(\mathbb{C}A\) and is such that A is the intersection of the sets \(\overline{u}^{1}(0)\) as u runs through H. Then every continuous real-valued function on X which is zero on A can be uniformly approximated by polynomials without constant terms in the functions of H.

Consider first the particular case in which A consists of a single point \(x_{0}\) . The hypotheses then imply that H separates the points of X; for, if \(x \neq x_{0}\) , then by hypothesis there is a function \(u \in H\) such that \(u(x) \neq 0 = u(x_{0})\) . Hence, for each \(\varepsilon > 0\) and each continuous real-valued function f on X such that \(f(x_{0}) = 0\) , there exists (Theorem 3) a polynomial g in the functions of H such that \(|f(x) - g(x)| \leqslant \varepsilon\) for all \(x \in X\) . In particular \(|g(x_{0})| \leqslant \varepsilon\) , so that

\[
\left| f (x) - \left(g (x) - g \left(x _ {0}\right)\right) \right| \leqslant 2 \varepsilon
\]

for all \(x \in \mathbf{X}\) ; and since \(g(x) - g(x_0)\) is a polynomial in the functions of \(\mathbf{H}\) with no constant term, the result is established in this case.

In the general case, consider the equivalence relation R on X whose classes are the set A and the sets \(\{x\}\) for \(x \notin A\) . The quotient space X/R is Hausdorff (Chapter I, § 8, no. 6, Proposition 15) and therefore compact. Let \(\varphi: X \to X/R\) be the canonical mapping. Every continuous real-valued function f on X which vanishes on A can be written in the form \(f = f_{1} \circ \varphi\) , where \(f_{1}\) is a continuous real-valued function on X/R which vanishes at the point \(x_{0}' = \varphi(A)\) . Applying the result already proved to the space X/R and the point \(x_{0}'\) , we obtain the final result.

